% Komprehensives Data Science & KI-Algorithmen Cheatsheet (K01-K14)
\scriptsize
\newgeometry{margin=.4cm}
\lstset{style=python, basicstyle=\ttfamily\tiny}
\pagestyle{empty}
\raggedcolumns
\tiny
\providecommand{\cheatSheetColor}{ForestGreen}

\newtcolorbox{cheatSheetBox}[1][]{
	enhanced jigsaw,
	colback=\cheatSheetColor!10,
	colframe=\cheatSheetColor,
	colbacktitle=\cheatSheetColor!30,
	fonttitle=\bfseries,
	parbox=false,
	attach boxed title to top text left={yshift=-5mm,yshifttext=-3.5mm},
	coltitle=black,
	#1}

\newcommand{\cheatSheetSubtitle}[2]{
	\textcolor{\cheatSheetColor}{
		\begin{centering}
			\subsection*{#1}
		\end{centering}
		\vspace{0.3em}
		\hrule
		#2
		\vspace{0.3em}
	}
}

\begin{landscape}
	\tiny
	\begin{multicols*}{4}
		\textcolor{\cheatSheetColor}{
			\begin{centering}
				\section*{KI \& Algorithmen -- Komprehensives Cheatsheet (K01--K14)}
			\end{centering}
			\hrule
			Das folgende Cheatsheet deckt alle Themen ab: Einführung, Datenvorbereitung, Algorithmen (Baum, RF, SVM, KNN, Regression, NN), Assoziation (Apriori), Filterung (Naive Bayes), Ranking (gewichtet, PageRank), Markov-Ketten und Korrelation/Kausalität.
			\vspace{0.2em}
		}

		% === K01: EINFÜHRUNG ===
		\cheatSheetSubtitle{K01: Einführung -- Algorithmen-Arten}{Unterscheidung nach Funktionsart und Einsatzbereich}
		
		\begin{cheatSheetBox}[title=Algorithmen nach Funktionsart]
			\textbf{1. Priorisierung} (Ranking): Gewichte, PageRank
			\textbf{2. Klassifikation} (Kategorien): Baum, RF, KNN, SVM, Markov
			\textbf{3. Assoziation} (Muster): Apriori
			\textbf{4. Filterung} (Zuweisung): Naive Bayes
			\textbf{5. Regression} (Zahlenwerte): Lin. Regression, NN
		\end{cheatSheetBox}

		\begin{cheatSheetBox}[title=Klassifikation vs. Regression]
			\textbf{Klassifikation}: Kategorie vorhersagen (diskret)
			\begin{lstlisting}
df = df[df['Note'] < 4]  # Filter nach Kategorie
			\end{lstlisting}
			\textbf{Regression}: Zahlenwert vorhersagen (kontinuierlich)
			\begin{lstlisting}
y = a0 + a1*x1 + a2*x2  # Lineare Gleichung
			\end{lstlisting}
		\end{cheatSheetBox}

		% === K02: DATENVORBEREITUNG ===
		\cheatSheetSubtitle{K02: Datenvorbereitung \& Evaluation}{Preprocessing und Metriken}

		\begin{cheatSheetBox}[title=Min-Max-Normalisierung]
			\[x' = \frac{x - x_{\min}}{x_{\max} - x_{\min}}\]
			Bringt alle Features auf $[0, 1]$ — essentiell für KNN, SVM
		\end{cheatSheetBox}

		\begin{cheatSheetBox}[title=Klassifikation-Metriken]
			\[
				\text{Accuracy} = \frac{TP + TN}{TP+TN+FP+FN}
			\]
			\[
				\text{Precision} = \frac{TP}{TP+FP}, \quad
				\text{Recall} = \frac{TP}{TP+FN}
			\]
			\textbf{TP}: Korrekt positiv, \textbf{FP}: Falsch positiv, \textbf{FN}: Falsch negativ
		\end{cheatSheetBox}

		\begin{cheatSheetBox}[title=Regression-Metriken]
			\[
				\text{MAE} = \frac{1}{n}\sum|y_i - \hat{y}_i|
			\]
			\[
				\text{RMSE} = \sqrt{\frac{1}{n}\sum(y_i - \hat{y}_i)^2}
			\]
			RMSE bestraft große Fehler stärker (Quadrierung)
		\end{cheatSheetBox}

		% === K03: ENTSCHEIDUNGSBAUM ===
		\cheatSheetSubtitle{K03: Entscheidungsbaum}{Stufenweise Aufteilung mit Gini-Koeffizient}

		\begin{cheatSheetBox}[title=Gini-Reinheitsmass]
			\[
				\text{Gini}(t) = 1 - \sum_{k=1}^{K} p_k^2
			\]
			Gini=0: Rein, Gini=0.5: Maximal unrein. Baum minimiert Gini iterativ.
		\end{cheatSheetBox}

		\begin{cheatSheetBox}[title=Overfitting-Prävention]
			\texttt{max\_depth}, \texttt{min\_samples\_split}, \texttt{min\_samples\_leaf} limitieren Tiefe und Komplexität
		\end{cheatSheetBox}

		% === K04: LINEARE REGRESSION ===
		\cheatSheetSubtitle{K04: Lineare Regression}{Vorhersage kontinuierlicher Werte}

		\begin{cheatSheetBox}[title=Einfache lineare Regression]
			\[y = mx + q\]
			Methode kleinster Quadrate:
			\[
				m = \frac{n\sum x_i y_i - \sum x_i \sum y_i}{n\sum x_i^2 - (\sum x_i)^2}
			\]
			\[
				q = \frac{\sum y_i - m\sum x_i}{n}
			\]
		\end{cheatSheetBox}

		\begin{cheatSheetBox}[title=Multiple lineare Regression]
			\[
				y = a_0 + a_1 x_1 + a_2 x_2 + \cdots + a_n x_n
			\]
			\begin{lstlisting}
model = smf.ols('Note ~ Schlaf + Lernaufwand', 
                 data=df).fit()
pred = model.predict({'Schlaf': [7], 'Lernaufwand': [3]})
			\end{lstlisting}
		\end{cheatSheetBox}

		% === K05: SVM ===
		\cheatSheetSubtitle{K05: SVM -- Support Vector Machine}{Optimale Trennebene mit maximalem Margin}

		\begin{cheatSheetBox}[title=SVM-Grundkonzept]
			Findet Hyperebene mit größtem Abstand (Margin) zwischen Klassen
			\begin{itemize}
				\item \textbf{Linearer Kernel}: Einfach, erklärbar
				\item \textbf{RBF-Kernel}: Flexibel, nichtlinear
			\end{itemize}
		\end{cheatSheetBox}

		\begin{cheatSheetBox}[title=Parameter C (Regularisierung)]
			$C$ groß: Enger Margin, Overfitting-Risiko
			$C$ klein: Breiter Margin, kann zu Underfitting führen
		\end{cheatSheetBox}

		% === K06: KNN ===
		\cheatSheetSubtitle{K06: k-Nearest Neighbors}{Mehrheitsabstimmung der k nächsten Nachbarn}

		\begin{cheatSheetBox}[title=Euklidische Distanz]
			\[
				d(a, b) = \sqrt{\sum_{i=1}^n (a_i - b_i)^2}
			\]
			\textbf{Kritisch}: Normalisierung essentiell! Größere Features dominieren sonst.
		\end{cheatSheetBox}

		\begin{cheatSheetBox}[title=Hyperparameter k]
			k=1: Sehr sensitiv, Overfitting
			k=10--30: Balance
			k groß: Glatte Grenzen, Underfitting-Risiko
		\end{cheatSheetBox}

		\columnbreak

		% === K07: NEURONALE NETZE ===
		\cheatSheetSubtitle{K07-K08: Neuronale Netze}{Schichtenweise Funktionsapproximation}

		\begin{cheatSheetBox}[title=Perceptron (Einzelneuron)]
			\[
				\hat{y} = f\left(\sum_{i=1}^{n} w_i x_i + b\right)
			\]
			$f$ = Aktivierungsfunktion (Sigmoid, ReLU, Softmax)
		\end{cheatSheetBox}

		\begin{cheatSheetBox}[title=Aktivierungsfunktionen]
			\textbf{Sigmoid}: $\sigma(z)=\frac{1}{1+e^{-z}}$ (0--1)
			\textbf{ReLU}: $\max(0,z)$ (0--$\infty$)
			\textbf{Softmax}: $\frac{e^{z_k}}{\sum_j e^{z_j}}$ (Wahrscheinlichkeiten)
		\end{cheatSheetBox}

		\begin{cheatSheetBox}[title=Backpropagation]
			Fehler rückwärts propagieren, Gewichte durch Gradientenabstieg anpassen
		\end{cheatSheetBox}

		\begin{cheatSheetBox}[title=NN-Grenzen]
			Black-Box (unerklärbar), hoher Datenbedarf, langsames Training, Overfitting-anfällig
		\end{cheatSheetBox}

		% === K09: MODELLVERGLEICH ===
		\cheatSheetSubtitle{K09: Modellvergleich-Übersicht}{Stärken, Grenzen, Einsatzbereiche}

		\begin{cheatSheetBox}[title=Algorithmen-Vergleich]
			\textbf{Baum}: Interpretierbar, Overfitting, kein Normalisierungsnötig
			\textbf{RF}: Robust, oft beste Accuracy, weniger erklärbar
			\textbf{SVM}: Effektiv bei klaren Grenzen, Parameterwahl komplex
			\textbf{KNN}: Einfach, keine Training, skalierungssensitiv
			\textbf{Lin. Reg.}: Schnell, linear, unerklärbar bei Abweichungen
			\textbf{NN}: Sehr mächtig, Black-Box, viele Daten nötig
		\end{cheatSheetBox}

		% === K10: TRANSFERAUFGABE ===
		\cheatSheetSubtitle{K10: Transferaufgabe}{Eigenes ML-Projekt-Checkliste}

		\begin{cheatSheetBox}[title=Schritte]
			1. Problem definieren (Klassifikation/Regression)
			2. Mindestens 50 gelabelte Zeilen sammeln
			3. Fehlende Werte, Normalisierung, 80/20-Split
			4. Zwei Modelle trainieren
			5. Metriken vergleichen, Visualisierung
			6. Bias-Risiken reflektieren
		\end{cheatSheetBox}

		% === K11: APRIORI ===
		\cheatSheetSubtitle{K11: Apriori -- Assoziationsregeln}{Häufige Itemsets und Regeln}{K11: Apriori}

		\begin{cheatSheetBox}[title=Zentrale Kennzahlen]
			\[
				\text{Support}(X \Rightarrow Y) = P(X \cup Y)
			\]
			\[
				\text{Confidence}(X \Rightarrow Y) = \frac{P(X \cup Y)}{P(X)}
			\]
			\[
				\text{Lift}(X \Rightarrow Y) = \frac{P(X \cup Y)}{P(X)P(Y)}
			\]
			Lift > 1: Positive Assoziation (stärker als Zufall)
		\end{cheatSheetBox}

		\begin{cheatSheetBox}[title=Apriori-Prinzip]
			Häufiges Itemset $\Rightarrow$ alle Teilmengen häufig
			Nicht-häufige Teilmenge $\Rightarrow$ keine häufige Obermenge
		\end{cheatSheetBox}

		% === K12: NAIVE BAYES ===
		\cheatSheetSubtitle{K12: Naive Bayes -- Filterung}{Wahrscheinlichkeitsbasierte Klassifikation}

		\begin{cheatSheetBox}[title=Bayes-Theorem]
			\[
				P(C\mid x) = \frac{P(x\mid C) \, P(C)}{P(x)}
			\]
			Naive Annahme: Bedingte Unabhängigkeit der Features
		\end{cheatSheetBox}

		\begin{cheatSheetBox}[title=Laplace-Glättung]
			\[
				P(w\mid C) = \frac{\text{count}(w,C) + \alpha}{\sum_v \text{count}(v,C) + \alpha|V|}
			\]
			Verhindert Wahrscheinlichkeit 0 für unbekannte Features
		\end{cheatSheetBox}

		% === K13: GEWICHTETES RANKING ===
		\cheatSheetSubtitle{K13: Gewichtetes Ranking}{Priorisierung durch Scoring}

		\begin{cheatSheetBox}[title=Scoring-Formel]
			\[
				\text{Score}(x) = \sum_{i=1}^n w_i z_i
			\]
			Typisch: $\sum w_i = 1$, Features normiert auf $[0,1]$
			\textbf{Warnung}: Gewichte sind Wertsetzungen, nicht objektiv!
		\end{cheatSheetBox}

		% === K14: PAGERANK ===
		\cheatSheetSubtitle{K14: PageRank -- Netzwerk-Ranking}{Wichtigkeit durch Backlinks}

		\begin{cheatSheetBox}[title=PageRank-Formel]
			\[
				PR(p) = \frac{1-d}{N} + d\sum_{q \to p}\frac{PR(q)}{\text{outdeg}(q)}
			\]
			$d$ = Dämpfung (typisch 0.85), $N$ = Seiten, outdeg = ausgehende Links
		\end{cheatSheetBox}

		\begin{cheatSheetBox}[title=Iterativer Prozess]
			Start: $PR = 1/N$ für alle. Dann iterativ aktualisieren bis konvergiert.
		\end{cheatSheetBox}

		\columnbreak

		% === MARKOV-KETTEN ===
		\cheatSheetSubtitle{Markov-Ketten}{Zustandsübergänge mit Wahrscheinlichkeiten}

		\begin{cheatSheetBox}[title=Übergangsmatrix (Wetter)]
			\[
				P = \begin{pmatrix} 0.8 & 0.2 \\ 0.4 & 0.6 \end{pmatrix}
			\]
			Zeilen: Aktueller Zustand (Sonne/Regen)
			Spalten: Nächster Zustand
		\end{cheatSheetBox}

		\begin{cheatSheetBox}[title=Vorhersage]
			\[
				\pi_1 = \pi_0 \times P
			\]
			Startvektor $\pi_0$ (z.B. $[1, 0]$ = Sonne) multipliziert mit Matrix
		\end{cheatSheetBox}

		\begin{cheatSheetBox}[title=Beispiel-Berechnung]
			Start: Sonne $\pi_0 = [1, 0]$
			Nach 1 Tag: $\pi_1 = [0.8, 0.2]$
			Nach 2 Tagen: $\pi_2 = [0.72, 0.28]$
		\end{cheatSheetBox}

		\begin{cheatSheetBox}[title=Python: Matrix und Vorhersage]
			\begin{lstlisting}
import numpy as np
P = np.array([[0.8, 0.2], [0.4, 0.6]])
pi0 = np.array([1, 0])
pi1 = pi0 @ P  # Nach 1 Tag
pi2 = pi1 @ P  # Nach 2 Tagen
			\end{lstlisting}
		\end{cheatSheetBox}

		\begin{cheatSheetBox}[title=Simulation mit Zufallszahlen]
			\begin{lstlisting}
import random
zustand = "Sonne"
for _ in range(20):
    if zustand == "Sonne":
        zustand = random.choices(
            ["Sonne", "Regen"], 
            weights=[0.8, 0.2]).pop()
    else:
        zustand = random.choices(
            ["Sonne", "Regen"], 
            weights=[0.4, 0.6]).pop()
    print(zustand)
			\end{lstlisting}
		\end{cheatSheetBox}

		% === KORRELATION & KAUSALITÄT ===
		\cheatSheetSubtitle{Korrelation \& Kausalität}{Unterscheidung, Irrtümer vermeiden}

		\begin{cheatSheetBox}[title=Korrelationskoeffizient]
			\[
				r = \frac{\sum(x_i - \bar{x})(y_i - \bar{y})}{\sqrt{\sum(x_i - \bar{x})^2 \cdot \sum(y_i - \bar{y})^2}}
			\]
			$r=1$: Perfekt positiv, $r=-1$: Perfekt negativ, $r=0$: Keine Korrelation
		\end{cheatSheetBox}

		\begin{cheatSheetBox}[title=Fünf Arten von Kausalität]
			1. \textbf{Koinzidenz}: Zufällige Korrelation (Mozzarella $\leftrightarrow$Doktortitel)
			2. \textbf{Zyklisch}: A$\rightarrow$B$\rightarrow$A Verstärkung
			3. \textbf{Gemeinsamer Grund}: C$\rightarrow$A, C$\rightarrow$B (Jahreszeit$\rightarrow$Verkauf+Erkältung)
			4. \textbf{Indirekt}: A$\rightarrow$C$\rightarrow$B (Sport$\rightarrow$Schlaf$\rightarrow$Noten)
			5. \textbf{Direkt}: A$\rightarrow$B
		\end{cheatSheetBox}

		\begin{cheatSheetBox}[title=Merksatz]
			\textbf{Korrelation $\neq$ Kausalität!} Immer auf Confounder prüfen.
		\end{cheatSheetBox}

		\columnbreak

		% === PRAKTISCHE PANDAS & PYTHON ===
		\cheatSheetSubtitle{Praktische Datenanalyse mit pandas}{Code-Snippets}

		\begin{cheatSheetBox}[title=CSV laden und inspizieren]
			\begin{lstlisting}
import pandas as pd
df = pd.read_csv('daten.csv')
df.head()       # Erste 5 Zeilen
df.info()       # Datentypen, Fehlwerte
df.describe()   # Statistiken
			\end{lstlisting}
		\end{cheatSheetBox}

		\begin{cheatSheetBox}[title=Spalten filtern und manipulieren]
			\begin{lstlisting}
spalte = df['Name']          # Eine Spalte
mehrere = df[['Name', 'Alter']]  # Mehrere
df_gefiltert = df[df['Alter'] > 30]  # Filter
df['BMI'] = df['Gewicht'] / (df['Grösse']/100)**2  # Neue Spalte
			\end{lstlisting}
		\end{cheatSheetBox}

		\begin{cheatSheetBox}[title=Statistiken und Aggregation]
			\begin{lstlisting}
df['Alter'].mean()       # Mittelwert
df['Alter'].sum()        # Summe
df['Alter'].std()        # Standardabweichung
df.groupby('Kategorie')['Wert'].mean()  # Gruppe aggregieren
df['Kategorie'].value_counts()  # Häufigkeiten
			\end{lstlisting}
		\end{cheatSheetBox}

		\begin{cheatSheetBox}[title=Visualisierung mit matplotlib]
			\begin{lstlisting}
import matplotlib.pyplot as plt
df['Spalte'].plot.hist(bins=20)  # Histogramm
df.plot.scatter(x='X', y='Y')  # Streudiagramm
df['Kategorie'].value_counts().plot.bar()  # Balken
df.plot.line(x='Zeit', y='Wert')  # Linie
			\end{lstlisting}
		\end{cheatSheetBox}

		% === SCIKIT-LEARN BEISPIELE ===
		\cheatSheetSubtitle{ML-Modelle mit scikit-learn}{Train-Fit-Predict Pattern}

		\begin{cheatSheetBox}[title=Train-Test-Split]
			\begin{lstlisting}
from sklearn.model_selection import train_test_split
X_train, X_test, y_train, y_test = train_test_split(
    X, y, test_size=0.2, random_state=42)
			\end{lstlisting}
		\end{cheatSheetBox}

		\begin{cheatSheetBox}[title=Entscheidungsbaum]
			\begin{lstlisting}
from sklearn.tree import DecisionTreeClassifier
model = DecisionTreeClassifier(max_depth=5)
model.fit(X_train, y_train)
y_pred = model.predict(X_test)
accuracy = model.score(X_test, y_test)
			\end{lstlisting}
		\end{cheatSheetBox}

		\begin{cheatSheetBox}[title=Random Forest]
			\begin{lstlisting}
from sklearn.ensemble import RandomForestClassifier
model = RandomForestClassifier(n_estimators=100, max_depth=8)
model.fit(X_train, y_train)
y_pred = model.predict(X_test)
print(model.feature_importances_)  # Feature-Gewichte
			\end{lstlisting}
		\end{cheatSheetBox}

		\begin{cheatSheetBox}[title=KNN]
			\begin{lstlisting}
from sklearn.neighbors import KNeighborsClassifier
model = KNeighborsClassifier(n_neighbors=5)
model.fit(X_train, y_train)
y_pred = model.predict(X_test)
			\end{lstlisting}
		\end{cheatSheetBox}

		\begin{cheatSheetBox}[title=SVM]
			\begin{lstlisting}
from sklearn.svm import SVC
model = SVC(C=1.0, kernel='rbf')
model.fit(X_train, y_train)
y_pred = model.predict(X_test)
			\end{lstlisting}
		\end{cheatSheetBox}

		\begin{cheatSheetBox}[title=Normalisierung]
			\begin{lstlisting}
from sklearn.preprocessing import MinMaxScaler
scaler = MinMaxScaler()
X_scaled = scaler.fit_transform(X_train)
X_test_scaled = scaler.transform(X_test)
			\end{lstlisting}
		\end{cheatSheetBox}

		% === EVALUATIONSFUNKTIONEN ===
		\cheatSheetSubtitle{Evaluation \& Metriken}{Genauigkeit überprüfen}

		\begin{cheatSheetBox}[title=Klassifikations-Metriken]
			\begin{lstlisting}
from sklearn.metrics import (accuracy_score, 
    precision_score, recall_score, f1_score)
accuracy = accuracy_score(y_test, y_pred)
precision = precision_score(y_test, y_pred)
recall = recall_score(y_test, y_pred)
f1 = f1_score(y_test, y_pred)
			\end{lstlisting}
		\end{cheatSheetBox}

		\begin{cheatSheetBox}[title=Regression-Metriken]
			\begin{lstlisting}
from sklearn.metrics import mean_absolute_error as mae
from sklearn.metrics import mean_squared_error as mse
mae_val = mae(y_test, y_pred)
mse_val = mse(y_test, y_pred)
rmse_val = np.sqrt(mse_val)
			\end{lstlisting}
		\end{cheatSheetBox}

		\begin{cheatSheetBox}[title=Konfusionsmatrix visualisieren]
			\begin{lstlisting}
from sklearn.metrics import confusion_matrix
import matplotlib.pyplot as plt
cm = confusion_matrix(y_test, y_pred)
plt.imshow(cm, cmap='Blues')
plt.xlabel('Vorhergesagt')
plt.ylabel('Tatsächlich')
			\end{lstlisting}
		\end{cheatSheetBox}

		\columnbreak

		% === BONUS: APRIORI CODE ===
		\cheatSheetSubtitle{Bonus: Apriori Implementation}{Assoziationsregeln finden}

		\begin{cheatSheetBox}[title=Apriori mit MLxtend]
			\begin{lstlisting}
from mlxtend.frequent_patterns import apriori
from mlxtend.frequent_patterns import association_rules
import pandas as pd

# Häufige Itemsets (Support > 0.1)
freq_itemsets = apriori(data, 
    min_support=0.1, use_colnames=True)

# Regeln mit Confidence > 0.5
rules = association_rules(freq_itemsets, 
    metric="confidence", min_threshold=0.5)
print(rules[['antecedents', 'consequents', 
             'support', 'confidence', 'lift']])
			\end{lstlisting}
		\end{cheatSheetBox}

		% === BONUS: NAIVE BAYES CODE ===
		\cheatSheetSubtitle{Bonus: Naive Bayes}{Text-Klassifikation}

		\begin{cheatSheetBox}[title=Naive Bayes Classifier]
			\begin{lstlisting}
from sklearn.naive_bayes import MultinomialNB
from sklearn.feature_extraction.text import TfidfVectorizer

vectorizer = TfidfVectorizer()
X_tfidf = vectorizer.fit_transform(texts)
model = MultinomialNB()
model.fit(X_tfidf, labels)
pred = model.predict(vectorizer.transform(new_texts))
			\end{lstlisting}
		\end{cheatSheetBox}

		% === QUICK REFERENCE TABLE ===
		\cheatSheetSubtitle{Schnellreferenz -- Wann welcher Algorithmus?}{Entscheidungshilfe}

		\begin{cheatSheetBox}[title=Algorithmus-Wahl nach Szenario]
			\textbf{Hohe Interpretierbarkeit nötig} $\rightarrow$ Entscheidungsbaum
			\textbf{Größte Genauigkeit gewünscht} $\rightarrow$ Random Forest
			\textbf{Wenig Daten, schnell} $\rightarrow$ KNN oder Lin. Reg.
			\textbf{Klare nichtlineare Grenzen} $\rightarrow$ SVM mit RBF-Kernel
			\textbf{Text-Klassifikation} $\rightarrow$ Naive Bayes
			\textbf{Assoziationen in Transaktionen} $\rightarrow$ Apriori
			\textbf{Zustandsübergänge} $\rightarrow$ Markov-Ketten
			\textbf{Komplexe nichtlineare Muster} $\rightarrow$ Neuronales Netz
			\textbf{Seiten-Ranking} $\rightarrow$ PageRank
		\end{cheatSheetBox}

		% === COMMON PITFALLS ===
		\cheatSheetSubtitle{Häufige Fehler vermeiden}{Best Practices}

		\begin{cheatSheetBox}[title=Fehler]
			1. \textbf{Datenleckage}: Test-Features in Training leaken
			2. \textbf{Keine Normalisierung}: KNN/SVM leiden unter großen Wertebereichen
			3. \textbf{Overfitting}: Zu komplexe Modelle, zu wenig Daten
			4. \textbf{Korrelation=Kausalität}: Confounder übersehen
			5. \textbf{Bias-Perpetuierung}: Historische Diskriminierung in Trainingsdaten
			6. \textbf{Zu wenig Evaluierung}: Nur Accuracy reicht nicht
		\end{cheatSheetBox}

		% === KEY FORMULAS ===
		\cheatSheetSubtitle{Zentrale Formeln Übersicht}{Referenz}

		\begin{cheatSheetBox}[title=Zentrale Formeln]
			\textbf{Lin. Regression}:
			$y = a_0 + a_1x_1 + \cdots + a_nx_n$

			\textbf{Gini}:
			$1 - \sum p_k^2$

			\textbf{Euklidische Distanz}:
			$\sqrt{\sum(a_i - b_i)^2}$

			\textbf{Bayes}:
			$P(C|x) = \frac{P(x|C)P(C)}{P(x)}$

			\textbf{PageRank}:
			$\frac{1-d}{N} + d\sum\frac{PR(q)}{\text{outdeg}(q)}$

			\textbf{Scoring}:
			$\sum w_i z_i$ (mit $\sum w_i=1$)
		\end{cheatSheetBox}

	\end{multicols*}
\end{landscape}
\clearpage
\restoregeometry